\chapter{Geometric Derivations of the Cosmic Energy Budget and Macroscopic Chronology}
\label{app:energy_and_chronology}
\addtocontents{toc}{\protect\setcounter{tocdepth}{0}}

\textbf{Abstract:} This appendix provides the strict mathematical derivations for the macroscopic geometric limits and the thermodynamic timeline of the universe. By modeling the active Spacetime Gravitational Floor (SGF) as a dual-phase visco-elastic Face-Centered Cubic (FCC) medium undergoing localized phase transitions, the exact volumetric capacity of the interstitial voids, the global percolation threshold for Catastrophic Decompression, and the absolute structural lifespan of the cosmic oscillator are derived as invariant mechanical boundaries.

\section{The Local Volumetric Yield Limit (Phase Condensation)}
As established in Part I, the active vacuum (State I) is an FCC lattice with an invariant geometric packing fraction of $\beta_{geo} \approx 0.74048$. Mass creation requires the localized kinematic excitation of the lattice, which reduces the effective structural tension to $\approx 69\%$. This localized mass-energy condenses specifically within the primary pockets of the interstitial condensate: the Octahedral voids.

To determine the absolute maximum mass the universe can locally sustain before the active vacuum geometrically exhausts and collapses into the frozen Simple Cubic core (State II), we calculate the precise volumetric fraction of these Octahedral voids relative to the FCC unit cell.

\begin{itemize}
    \item \textbf{Deriving the Octahedral Void Radius ($r_O$):} 
    In an FCC lattice formed by nodes of radius $R$, the nodes touch along the face diagonal. The face diagonal is therefore $4R$. By the Pythagorean theorem, the edge length of the unit cell ($a$) is:
    $$a = \sqrt{\frac{(4R)^2}{2}} = 2\sqrt{2} R$$
    An Octahedral void is located perfectly at the midpoint of each edge of the unit cell. Looking along this edge, the distance from the center of a corner node to the center of the edge void is exactly half the edge length ($a/2$). Therefore, the radius of the node ($R$) plus the inscribed radius of the void ($r_O$) must equal $a/2$:
    $$R + r_O = \frac{a}{2} = \frac{2\sqrt{2} R}{2} = \sqrt{2} R$$
    Isolating the void radius yields:
    $$r_O = \sqrt{2} R - R = (\sqrt{2}-1)R$$

    \item \textbf{Deriving the Volumetric Fraction:} 
    The total volume of the unit cell is $V_{cell} = a^3 = 16\sqrt{2} R^3$. Because there are exactly 4 total Octahedral voids per FCC unit cell (1 in the center, plus 12 shared quarter-voids on the edges), the maximum available void volume ($V_{oct}$) is:
    $$V_{oct} = 4 \left( \frac{4}{3} \pi r_O^3 \right) = \frac{16\pi}{3}(\sqrt{2}-1)^3 R^3$$
    Dividing this specific void volume by the total cell volume yields the Octahedral Volumetric Fraction ($\beta_{oct}$):
    $$\beta_{oct} = \frac{V_{oct}}{V_{cell}} = \frac{\frac{16\pi}{3}(\sqrt{2}-1)^3 R^3}{16\sqrt{2} R^3} = \frac{\pi(\sqrt{2}-1)^3}{3\sqrt{2}} \approx 0.0526$$
\end{itemize}

The local volumetric capacity for macroscopic mass formation ($\beta_b$) is strictly bounded at 5.26\% ($\beta_b \leq \beta_{oct}$). Condensation beyond this limit mathematically exhausts the interstitial capacity, triggering the $\rho_{QG}$ phase transition (Black Hole formation).

\section{The Global Topological Limit (Percolation Theory and The Big Rupture)}

While the volumetric limit bounds the local structural capacity for mass condensation, Percolation Theory bounds the ultimate global stability of the macroscopic cosmic reservoir. A crystalline tension field completely ruptures only when its topological defects---the frozen State II Black Hole cores left behind after local phase collapses---mechanically connect to form a contiguous, macroscopic fault line. 

\begin{itemize}
    \item \textbf{The Interwoven Sub-Lattice:} In standard crystallography, the spatial coordinates of the primary interstitial voids within an FCC matrix arrange themselves into a secondary, interwoven geometric pattern known as a sub-lattice. This is not a physical barrier, but a strict coordinate map of the available interstitial condensate capacity. State II structural defects (hyper-dense SC nodes) form exclusively at these specific geometric sites.
    
    \item \textbf{Site Percolation Threshold ($p_c$):} The established mathematical constant for site percolation (the threshold for global structural connectivity) in an FCC lattice is $p_c \approx 0.199$.
    
    \item \textbf{Total Interstitial Fraction ($\beta_{IC}$):} As established in Section 6, the total available interstitial condensate capacity in the active vacuum is $\beta_{IC} = 1 - \beta_{geo} \approx 0.2595$.
\end{itemize}

To find the absolute global threshold for defect density ($\beta_{\text{rup}}$) before the cosmic reservoir suffers Catastrophic Decompression, we multiply the sub-lattice percolation threshold by the total interstitial capacity:
$$\beta_{\text{rup}} = p_c \times \beta_{IC} = 0.199 \times 0.2595 \approx 0.0516$$
Because the global percolation limit ($5.16\%$) is mathematically lower than the local void limit ($5.26\%$), the universe is structurally guaranteed to shatter globally before it can ever exhaust its local mass-generation capacity.

\subsection{Derivation of the Deflation Parameters}
\label{subsec:deflation_parameters}

Once the absolute threshold of $5.16\%$ is established, we can derive the strict physical boundary limits of the deflating cosmic reservoir at the exact moment of this rupture. As active Face-Centered Cubic (FCC) vacuum is consumed and converted into zero-void SC cores, the macroscopic reservoir must physically deflate to maintain the contiguous tension of the field.

\begin{itemize}
    \item \textbf{Global Rupture Trigger Mass ($M_{\text{trg}}$):} The absolute mass threshold at which Catastrophic Decompression occurs. The rupture triggers strictly when the Aggregate Percolating Mass ($M_p$) reaches $5.16\%$ of the total macroscopic reservoir mass ($M_{\text{Res}}$).
    $$M_{\text{trg}} = 0.0516 \cdot (1.9851 \times 10^{79} \text{ kg}) \approx 1.0243 \times 10^{78} \text{ kg}$$
    
    \item \textbf{Volumetric Deflation Limit ($V_{\text{rup}}$):} The active FCC volume shrinks by exactly $5.16\%$, leaving $94.84\%$ of the initial volume intact at the moment of structural failure.
    $$V_{\text{rup}} = 0.9484 \cdot V_{\text{Res}} \approx 9.409 \times 10^{104} \text{ m}^3$$
    
    \item \textbf{Deflated Rupture Radius ($R_{\text{rup}}$):} The physical contraction of the bounding reservoir boundary at the moment the fault line snaps is derived from the cube root of the final volume. 
    $$R_{\text{rup}} = R_{\text{Res}} \times (0.9484)^{1/3} \approx 6.077 \times 10^{34} \text{ m}$$
\end{itemize}
The boundary of the Geometric Cosmic Horizon physically shrinks inward by $\approx 1.09 \times 10^{33} \text{ m}$ prior to catastrophic failure.

\subsection{The Spanning Cluster and Topological Continuity}

In statistical mechanics, the percolation threshold ($p_c$) marks the geometric phase transition from isolated local clusters to a single macroscopic ``Spanning Cluster.'' It may seem counterintuitive that a volumetric fill of only $5.16\%$ is mathematically sufficient to form a continuous physical chain stretching across the entire deflated diameter of $6.077 \times 10^{34} \text{ m}$. To prove this geometric capacity, we analyze the available linear span of the defect network.

Let the total rupture volume of the universe be $V_{\text{rup}} = \frac{4}{3}\pi R_{\text{rup}}^3$. The total volume of all SC defects is exactly $V_{\text{def}} = 0.0516 \cdot V_{\text{rup}}$.

Assuming this defect volume is distributed into $N$ identical macroscopic Black Holes, each with a localized radius $r$, the total number of Black Holes is the total defect volume divided by the individual volume:
$$N = \frac{V_{\text{def}}}{\frac{4}{3}\pi r^3}$$

The total maximum linear length ($L_{\text{tot}}$) available to construct the fault line, if these $N$ SC cores were aligned shoulder-to-shoulder, is their total count multiplied by their diameter ($2r$):
$$L_{\text{tot}} = N(2r) = \left( \frac{V_{\text{def}}}{\frac{4}{3}\pi r^3} \right)(2r) = \frac{3 V_{\text{def}}}{2\pi r^2}$$

\textbf{Geometric Proof:} Because the individual defect radius ($r$) is squared in the denominator, the total linear length of the spanning material scales inversely with the square of the defect size ($L_{\text{tot}} \propto r^{-2}$). By distributing $5.16\%$ of the universe's volume into stellar-mass SC cores, the resulting linear geometry is immensely vast. It provides orders of magnitude more physical length than is required to span the reservoir. Consequently, the distributed cores possess more than enough structural material to form a highly tortuous, unbroken fractal web.

\subsection{The Phase Trigger Unification}

This topological limit provides the exact mechanical cause for the Big Rupture criteria. As the universe ages, it accumulates frozen State II cores. When this distributed network of rigid defects reaches $5.16\%$ of the reservoir's capacity, the network globally percolates. 

This mechanically contiguous fault line physically shatters the structural integrity of the active vacuum, forcing the containment pressure gradient to vanish ($\Delta P \to 0$). Without this stabilizing pressure, the hyper-compressed fault line violently yields, triggering the martensitic snap-back (Catastrophic Decompression) that restores the macroscopic Face-Centered Cubic (FCC) geometry for the next cosmic cycle. The cosmic energy budget is therefore governed strictly by the invariant geometric packing limit of the active vacuum.

\section{The Chronology of the Global Phase Avalanche}

Because the topological limit ($\beta_{\text{rup}} \approx 0.0516$) dictates the ultimate structural failure of the active vacuum, it simultaneously governs the thermodynamic timeline of the macroscopic oscillator.

\subsection{The Global Phase Avalanche Equation}

To isolate the absolute thermodynamic lifespan of the universe, we mathematically evaluate the macroscopic phase clock by tracking the exponential growth required to scale the fundamental Micro-Black Hole Seed ($M_{D,G}$) to the global defect limit. 

Because Black Holes act as the ultimate thermodynamic sinks of the universe, this global mass-energy growth is strictly driven by the fundamental energy dissipation rate of the contiguous FCC field ($\Gamma_D$). We define the Global Phase Avalanche Equation as:
$$M_p(t) = M_{D,G} \cdot e^{\Gamma_D t}$$

\textbf{Variable Definitions:}
\begin{itemize}
    \item \textbf{$M_p(t)$ (The Aggregate Percolating Mass):} The total combined mass-energy segregated onto the static State II phase boundaries by the entire global network of Black Holes at time $t$. This equation models the macroscopic thermodynamic accumulation of the entire reservoir leading strictly toward the global percolation threshold.
    \item \textbf{$M_{D,G}$ (The Micro-Black Hole Seed):} $\approx 10.625 \text{ kg}$. The absolute minimum kinematic threshold required to fracture a localized void and nucleate a State II core.
    \item \textbf{$\Gamma_D$ (The Vacuum Damping Rate):} The visco-elastic compounding rate driving the exponential mass accumulation. The inverse of this constant ($1/\Gamma_D$) defines the fundamental visco-elastic decay time.
\end{itemize}

\textbf{The Initialization Boundary ($t \approx 0$)}

Because State II cores are strictly geometric lattice collapses rather than atomic structures, their formation does not require the universe to cool. In the immediate, highly chaotic aftermath of Catastrophic Decompression ($t \approx 0$), overlapping extreme acoustic shockwaves instantly push localized regions of the active vacuum beyond their kinematic saturation limit ($E_{\text{max}}$). 

This triggers the immediate formation of the first generation of State II cores (Primordial Black Holes) at $t \approx 0$, preceding the formation of stable atomic structures by 386,000 years (\hyperref[app:lattice_annealing]{Appendix~\ref*{app:lattice_annealing}}). Therefore, the exponential avalanche clock initiates at the exact moment the active vacuum is restored, functionally predating both atomic recombination and mature galactic evolution.

\subsection{Integration of the Macroscopic Lifespan ($T_{\text{cyc}}$)}

As derived in \hyperref[sec:reservoir_mass]{Section~\ref*{sec:reservoir_mass}}, the exact mass of the total cosmic reservoir is $M_{\text{Res}} \approx 1.9851 \times 10^{79} \text{ kg}$. The macroscopic cycle terminates when the defect network achieves the 19.9 percent site percolation threshold, requiring the conversion of exactly 5.16 percent of this total reservoir into rigid cores.

As derived in \hyperref[subsec:deflation_parameters]{Section~\ref*{subsec:deflation_parameters}}, the Global Rupture Trigger ($M_{\text{trg}}$) occurs at this precise topological limit ($1.0243 \times 10^{78} \text{ kg}$). Integrating the phase avalanche from the fundamental nucleation seed ($M_{D,G}$) to this rupture limit yields the total oscillation time ($T_{\text{cyc}}$):
$$1.0243 \times 10^{78} \text{ kg} = 10.625 \text{ kg} \cdot e^{\Gamma_D T_{\text{cyc}}}$$
$$e^{\Gamma_D T_{\text{cyc}}} \approx 9.6405 \times 10^{76}$$
$$\Gamma_D T_{\text{cyc}} = \ln(9.6405 \times 10^{76}) \approx 177.26$$

Multiplying by the visco-elastic decay time ($1/\Gamma_D \approx 13.39 \text{ billion years}$) establishes the absolute structural lifespan:
$$T_{\text{cyc}} = 177.26 \cdot 13.39 \text{ Billion Years} \approx 2.374 \text{ Trillion Years}$$

\subsection{Integration of the Current Epoch ($t_{\text{cur}}$)}

To calculate the current elapsed time ($t_{\text{cur}}$), we evaluate the equation against the aggregate mass of the currently existing State II defect network ($M_{\text{cur}}$).

Current astrophysical census data bounds the inertial mass of all Black Holes within our causal horizon ($R_U$) at $M_{\text{obs}} \approx 1.000 \times 10^{50} \text{ kg}$. This observable inertia strictly represents the mass-energy already harvested from the active vacuum and frozen into State II defects locally. 

To determine the current global aggregate mass, we scale this local harvested volume across the entire structural reservoir using the exact volumetric multiplier ($1.1633 \times 10^{26}$) established in \hyperref[sec:reservoir_mass]{Section~\ref*{sec:reservoir_mass}}:
$$M_{\text{cur}} = (1.000 \times 10^{50} \text{ kg}) \cdot (1.1633 \times 10^{26}) = 1.1633 \times 10^{76} \text{ kg}$$

Integrating the clock backward from this current aggregate mass establishes the present epoch:
$$1.1633 \times 10^{76} \text{ kg} = 10.625 \text{ kg} \cdot e^{\Gamma_D t_{\text{cur}}}$$
$$e^{\Gamma_D t_{\text{cur}}} \approx 1.0948 \times 10^{75}$$
$$\Gamma_D t_{\text{cur}} = \ln(1.0948 \times 10^{75}) \approx 172.88$$
$$t_{\text{cur}} = 172.88 \cdot 13.39 \text{ Billion Years} \approx 2.315 \text{ Trillion Years}$$

\subsection{The Expiration Timer and Remaining Lifespan}

To find the exact differential time remaining ($\Delta t = T_{\text{cyc}} - t_{\text{cur}}$) before Catastrophic Decompression, we evaluate the ratio of the global rupture limit to the current aggregate mass. Dividing the two phase states cancels the fundamental nucleation seed ($M_{D,G}$) and isolates the temporal difference:
$$\frac{M_{\text{trg}}}{M_{\text{cur}}} = \frac{M_{D,G} \cdot e^{\Gamma_D T_{\text{cyc}}}}{M_{D,G} \cdot e^{\Gamma_D t_{\text{cur}}}} = e^{\Gamma_D (T_{\text{cyc}} - t_{\text{cur}})} = e^{\Gamma_D \Delta t}$$

Taking the natural logarithm of both sides allows us to strictly bound the remaining lifespan:
$$\ln\left(\frac{M_{\text{trg}}}{M_{\text{cur}}}\right) = \Gamma_D \Delta t$$
$$\Delta t = \frac{1}{\Gamma_D} \ln\left(\frac{M_{\text{trg}}}{M_{\text{cur}}}\right)$$

Evaluating the exact mass ratio:
$$\frac{M_{\text{trg}}}{M_{\text{cur}}} = \frac{1.0243 \times 10^{78} \text{ kg}}{1.1633 \times 10^{76} \text{ kg}} \approx 88.05$$

Substituting this ratio into the differential time equation reveals the exact physical time remaining on the macroscopic clock:
$$\Delta t = \frac{1}{\Gamma_D} \ln(88.05) \approx 4.478 \cdot \frac{1}{\Gamma_D}$$

This mathematically guarantees that the active vacuum has approximately $4.478$ visco-elastic decay cycles ($e$-foldings) remaining before the topological defect network achieves global percolation. By multiplying by the decay constant, the remaining lifespan of the macroscopic universe is mechanically determined to be precisely $59.96$ billion years.

\subsection{The Macroscopic Phase Symmetry (Optical and Thermodynamic Equivalence)}

As established in the main text \hyperref[sec:causal_distance]{Section \ref*{sec:causal_distance}}, the absolute structural lifespan of the active vacuum ($T_{\text{cyc}}$) defines the Maximum Causal Length ($L_{\text{max}} = T_{\text{cyc}} \cdot c_E$). While a transverse wave mathematically thermalizes into the CMB noise floor earlier at the Extinction Length ($L_{\text{ext}}$), evaluating the decay equation strictly at the absolute kinematic boundary ($L_{\text{max}}$) reveals the exact analytical symmetry of the macroscopic oscillator.

The visco-elastic decay of a maximum-energy excitation ($E_{\text{max}}$) evaluated at the absolute causal boundary is:
$$E(L_{\text{max}}) = E_{\text{max}} \cdot e^{-L_{\text{max}}/R_U}$$

Because $L_{\text{max}} = c_E T_{\text{cyc}}$ and the viscous horizon is $R_U = c_E / \Gamma_D$, we substitute these kinematics into the exponent, isolating the exact dimensionless phase constant ($\Gamma_D T_{\text{cyc}}$) derived in the temporal timeline:
$$\frac{L_{\text{max}}}{R_U} = \frac{c_E T_{\text{cyc}}}{c_E / \Gamma_D} = \Gamma_D T_{\text{cyc}}$$
$$\frac{E(L_{\text{max}})}{E_{\text{max}}} = e^{-\Gamma_D T_{\text{cyc}}}$$

Simultaneously, the Global Phase Avalanche equation dictates that the accumulation of State II defect mass to the rupture threshold ($M_{\text{trg}}$) from the localized nucleation seed ($M_{D,G}$) is defined by the exact reciprocal of this phase constant:
$$\frac{M_{\text{trg}}}{M_{D,G}} = e^{+\Gamma_D T_{\text{cyc}}}$$
$$e^{-\Gamma_D T_{\text{cyc}}} = \frac{M_{D,G}}{M_{\text{trg}}}$$

By substituting the inverted mass ratio directly into the energy fraction, the exponential decay constant is completely eliminated, yielding a pure, invariant scaling identity:
$$\frac{E(L_{\text{max}})}{E_{\text{max}}} = \frac{M_{D,G}}{M_{\text{trg}}}$$

Isolating the final absolute energy state yields:
$$E(L_{\text{max}}) = E_{\text{max}} \cdot \left( \frac{M_{D,G}}{M_{\text{trg}}} \right)$$

Evaluating this identity using the established framework constants ($E_{\text{max}} = 25.18 \text{ J}$, $M_{D,G} = 10.625 \text{ kg}$, $M_{\text{trg}} = 1.0243 \times 10^{78} \text{ kg}$):
$$E(L_{\text{max}}) = 25.18 \text{ J} \cdot \left( \frac{10.625 \text{ kg}}{1.0243 \times 10^{78} \text{ kg}} \right) \approx 2.61 \times 10^{-76} \text{ Joules}$$

This analytical identity proves the scale-invariant geometric closure of Tension Field Theory. The fractional energy remaining in the universe's most powerful excitation exactly at the expiration of the macroscopic clock is geometrically identical to the mass fraction of a single microscopic Black Hole compared to the entire ruptured cosmic reservoir. The energy decay and the mass accumulation are exact mathematical inverses.
